\documentclass[11pt,a4paper]{article}
\usepackage[T1]{fontenc}
\usepackage[utf8]{inputenc}
\usepackage{amsmath,amssymb,amsthm,newtxtext,newtxmath,microtype}
\usepackage[margin=23mm,headheight=14pt,headsep=7mm,footskip=11mm]{geometry}
\usepackage{graphicx,booktabs,array,tabularx,enumitem,fancyhdr,caption,float,xcolor}
\usepackage[numbers,sort&compress]{natbib}
\usepackage[colorlinks=true,allcolors=blue!45!black,breaklinks=true]{hyperref}
\definecolor{ink}{HTML}{253640}\definecolor{accent}{HTML}{176B9B}
\pagestyle{fancy}\fancyhf{}\fancyhead[L]{\small Geometric root calculators}\fancyhead[R]{\small V22 / Research preprint}\fancyfoot[C]{\small\thepage}
\setlength{\parindent}{0pt}\setlength{\parskip}{5pt plus 1pt}\setlength{\emergencystretch}{3em}
\renewcommand{\arraystretch}{1.15}\captionsetup{font=small,labelfont=bf}
\setlist{noitemsep,topsep=3pt,leftmargin=*}\raggedbottom
\newtheorem{theorem}{Theorem}[section]\newtheorem{proposition}[theorem]{Proposition}
\newcommand{\map}[1]{\mathcal T_{\mathrm{#1}}}
\newcommand{\hub}[1]{\mathcal U_{#1}}
\newcommand{\fig}[3]{\begin{figure}[H]\centering\includegraphics[width=#2\linewidth]{figures/#1.pdf}\caption{#3}\end{figure}}
\hypersetup{pdftitle={Geometric Root Calculators: Fixed Circles, Global Convergence, and Analog Realization},pdfauthor={Ivan Besevic}}
\begin{document}\thispagestyle{empty}
{\small\sffamily\color{accent} CONSTRUCTIVE GEOMETRY / ITERATIVE COMPUTATION / ANALOG REALIZATION}\par
\vspace{8pt}
{\fontsize{24}{27}\selectfont\bfseries\color{ink} Geometric Root Calculators\par}
\vspace{5pt}{\large Fixed circles, global convergence, and analog realization}\par
\vspace{9pt}Ivan Besevic\footnote{Computational assistance: AI was used for derivations, proof development, figures and manuscript preparation. The author is responsible for the claims and their attribution. Section~\ref{sec:scope} distinguishes exact proofs, numerical checks and exploratory displays.}\quad$\cdot$\quad Independent research\par
{\small Version 22\quad$\cdot$\quad19 September 2026}\par
\url{https://github.com/ivan-fr/pandrosion}
\vspace{8pt}\hrule\vspace{8pt}
\textbf{Abstract.}
We construct positive-root iterations using lines and circles, with an explicit count of mobile traces. A prepared fixed circle solves an inverse Pad\'e model and gives a fifth-order update: eight joins for the cubic example, or $2p+2$ with a linear power stage. For every integer $p\ge3$, we prove that the selected real branch strictly reduces logarithmic error. Iteration therefore converges from every point in its real band. A centered AD step handles starts outside that band and exceptional readouts, giving a global algorithm whose control cost is counted separately. Four accelerated Pandrosion constructions and the linear MA reference provide geometric comparisons. We then derive a centered quadratic equation for a prospective analog cell and relate it to an existing order-three AD prototype. Lean certificates, executable constructions and accuracy-versus-work tables document the results.

\section{Why construct an iteration?}
A numerical algorithm names operations. A physical calculator must realize them: a ratio may be an alignment, a product a similar-triangle construction, and a square root a circle crossing. Two formulas with the same local order can therefore have different realization costs. This paper asks which geometric arrangements make a root iteration economical and how that arrangement can guide an analog implementation.

The intended setting is constructive computation: exact incidence geometry supplies a design and a correctness argument; a count of mobile operations describes the protocol. It also exposes which quantities must be stored, copied or solved for in a circuit. Drawing cost and electrical delay are different quantities, but the decomposition into reusable operations connects them.

Our principal example is a fixed circle that solves for a correction rather than merely transferring a length. Three joins after the power stage complete a fifth-order update. The circle converges throughout its real branch; a globally defined AD fallback handles entry and geometric exceptions. The decentered arc gives another choice: global monotone convergence with a moving radius. The contribution is these explicit realizations and their analysis, rather than the classical Newton, Halley or Pad\'e families. Section~\ref{sec:scope} collects the proof and implementation boundaries; the main text develops the constructions themselves.

\clearpage
\section{Coordinates as reusable computing elements}\label{sec:rails}
Let $X>0$, let $p\ge3$ be an integer, and write
\[
 a=X^{-1/p},\qquad t=Xs^p,\qquad s^+=s\Phi(t),\qquad s>0.
\]
The state approaches the reciprocal root $a$; inversion supplies the desired $X^{1/p}$. We use $e=\log(s/a)$ for local error. Points are capital letters, scalar iteration maps are $\map{label}$, and polynomial coefficients use Greek letters. Thus a point and a correction function never share the same symbol.

In a rectangle with vertices $C=(0,0)$, $B=(W,0)$, $O=(0,H)$ and $A=(W,H)$, define two extended rails and a family of top centers:
\begin{equation}\label{eq:rails}
 R(q)=(W,H(1-q)),\quad L(q)=(0,H(1-q)),\quad
 \hub{f}=\left(\frac{Wf}{f-1},H\right).
\end{equation}
A line through the prepared center performs multiplication in one direction and division in the other:
\begin{equation}\label{eq:railops}
 \hub{f}R(q)\cap OC=L(fq),\qquad
 \hub{f}L(b)\cap AB=R(b/f).
\end{equation}
For example, substituting the first two points into a line equation and setting $x=0$ gives ordinate $H(1-fq)$. This is the whole arithmetic content of the rail device.

\subsection{The power stage}
Start at $P=R(s)$. The join $CP$ meets the top line at $V=\hub{1/s}$. A leftward projection through $\hub{1/2}$ followed by a return through $V$ changes a rail value $q$ to $sq/2$. After $p-1$ pairs, the endpoint is $Q=R(s^p/2^{p-1})$. Including the initial join, the stage uses $2p-1$ lines. Its normalization factor is part of the design, not discarded arithmetic.

More generally, choosing the leftward factors $f_1,\ldots,f_{p-1}$ gives an endpoint with value $s^p\prod_i f_i$. This allows the power stage to feed a device with its preferred residual scale. We use exactly this freedom for the fixed-circle calculator.


\textbf{Power protocols, one state encoding.} The later rectangle drawings use a different power stage: similar triangles and parallels produce $E=R(s^p)$ directly. They do not change the meaning of $R(s)$. The rail telescope instead produces a scaled power; its scale is absorbed into the prepared readout. Their costs must be kept separate.
\begin{center}\small
\begin{tabularx}{\linewidth}{@{}lXX@{}}\toprule
Power stage & Geometric mechanism & Output and use\\\midrule
Rail telescope & $2p-1$ joins; repeated projection & Scaled $s^p$; pencils and fixed circle\\
Triangle chain & Horizontals and similar-triangle parallels & $R(s^p)$; original MA/AK/AD drawings\\
Binary products & Repeated squaring and multiplication & $R(s^p)$; large degrees and safe controllers\\\bottomrule
\end{tabularx}\end{center}

\subsection{A precise notion of work}
We count a join as $J$, a parallel-construction macro as $P_{\parallel}$, and a mobile compass arc as $C_{\mathrm{arc}}$. Intersections of already drawn objects, fixed preparation and final inversion are excluded. When expanded into ruler-and-compass traces, our parallel macro costs $2J+3C_{\mathrm{arc}}$. These conventions are specified before any cost comparison.

This question belongs to the constructive tradition of Lemoine's geom\'etrographie and d'Ocagne's nomography \cite{lemoine1902,docagne1925}. The Poncelet--Steiner theorem supplies a broad existence result for straightedge constructions with a given circle and its center \cite{steiner1833}. Our question is more specific: an explicit repeated protocol, its branch selection, and its mobile cost. The name Pandrosion refers to the addressee of the opening of Book III of Pappus' \emph{Collection} \cite{pappusLittle2023}. Our rectangle follows a modern reconstruction \cite{launayPandrosion}; the accelerated supports are modern constructions, not attributed to the ancient text.

\section{A prepared circle as a quadratic solver}\label{sec:circle}
Suppose the desired correction factor is $v\approx1$. Instead of evaluating a rational approximation to $t^{-1/p}$, approximate the inverse relation $t=v^{-p}$ by a quadratic rational model and solve it exactly. Put
\begin{align}
 \alpha_p&=(p+1)(p+2),&\beta_p&=p^2-4,&\gamma_p&=(p-1)(p-2),\nonumber\\
 N(v)&=\gamma_pv^2-2\beta_pv+\alpha_p,&
 D(v)&=\alpha_pv^2-2\beta_pv+\gamma_p.\label{eq:model}
\end{align}
The model is $\mathcal R(v)=N(v)/D(v)$, the $[2/2]$ Pad\'e approximant of $v^{-p}$ at one. Its inverse branch $v(1)=1$ defines $\map{circle}(s)=sv(Xs^p)$. A circle can supply the quadratic root in a single crossing.

\subsection{The eight-line cubic example}
Use $W=2$, $H=4$, $p=3$ and $X=2$. Prepare a circle through $B$ with center $O_\Gamma$, a projector $Z$, and constants
\[
 O_\Gamma=(-152/113,126/113),\quad Z=(478/113,396/113),
 \qquad \rho=\tfrac12,\quad k=\tfrac5{113}.
\]
The first five joins form the rail power stage with left factors $\rho$ and $\sigma=20/113$, retaining $L_1=L(\rho s)$ and ending at $T=R(kXs^3)$. The last three joins are the main device:
\begin{enumerate}
\item Draw $ZT$ and select its prescribed intersection $G$ with the fixed circle.
\item Draw $BG$ to the top line, obtaining $U=\hub{\rho/v}$.
\item Draw $UL_1$ back to the right rail. Its value is $(\rho s)/(\rho/v)=sv$.
\end{enumerate}
\fig{fixed_circle}{1}{The fixed-circle device at $s=3/4$. Left: the sixth join supplies $G$, and the seventh supplies the readout center $U$. Right: an independently enlarged view shows the eighth join and distinguishes $P$ from $P^+$. The circle is drawn once, before iteration. Equal horizontal and vertical scales are used within each view.}
The computed next state is $0.793700551706633\ldots$, with relative error about $3.24\times10^{-8}$. The circle is doing arithmetic: it replaces a longer rational readout by the solution of a quadratic model.

\clearpage
\subsection{Why the inverse correction has order five}
At $v=1$, $N=D=12$. The polynomial certificate for local order is
\begin{equation}\label{eq:cert}
 pND+v(N'D-ND')=p\alpha_p\gamma_p(v-1)^4.
\end{equation}
Moreover $\gamma_pN=(\gamma_pv-\beta_p)^2+3(p^2-4)$ and $\alpha_pD=(\alpha_pv-\beta_p)^2+3(p^2-4)$, so the logarithms used below are well defined for $p>2$.

\begin{theorem}[Local inverse-circle order]
For real $p>2$, the inverse branch through $v=1$ exists near $t=1$ and satisfies
\begin{equation}\label{eq:circleorder}
 e^+=-\frac{(p^2-1)(p^2-4)}{720}e^5+O(e^7).
\end{equation}
It therefore alternates sides locally and has exact order five.
\end{theorem}
\begin{proof}
Set $v=1+u$, $E(u)=\log\mathcal R(1+u)/p$ and $K(u)=E(u)+\log(1+u)$. Equation~\eqref{eq:cert} gives
\[
 K'(u)=\frac{\alpha_p\gamma_pu^4}{(1+u)N(1+u)D(1+u)}.
\]
Thus $K(u)/u^5\to\alpha_p\gamma_p/720$, while $E(0)=0$ and $E'(0)=-1$. The inverse function theorem identifies $E(u)$ with the input logarithmic error. Dividing the limits gives the coefficient in~\eqref{eq:circleorder}. Analyticity and $\mathcal R(1/v)=1/\mathcal R(v)$ make the logarithmic update odd, giving the stated remainder. The coefficient is nonzero for $p>2$.
\end{proof}
The same identity proves the Pad\'e contact by integrating $\frac{d}{dv}\log(v^p\mathcal R(v))$. Model-based iteration and quadratic approximants have classical precedents \cite{traub1961,traub1964,shafer1974}.

\subsection{Which circle intersection?}
The discriminant and the descending branch $\mathcal R'(v)<0$ are
\begin{align*}
 \Delta_p(t)&=(42p^2-24)t-3(p^2-4)(1+t^2),\\
 \tau_\pm&=\frac{7p^2-4\pm4p\sqrt{3(p^2-1)}}{p^2-4},\qquad \tau_-\tau_+=1,\\
 v_\pm&=\frac{p^2+2\pm\sqrt{12(p^2-1)}}{p^2-4}.
\end{align*}
The transverse real interval is $\tau_-<t<\tau_+$; select its inverse with $v_-<v<v_+$.
Positivity alone does not select this branch: the figure exhibits two positive roots at the same residual. At $\Delta_p(t)=0$ they merge tangentially. Theorem~\ref{thm:branch} proves contraction in the open band; AD handles tangencies and starts outside it.
\fig{branch_fold}{.82}{The fold of $\mathcal R$ for $p=3$. Blue marks the descending branch through $(1,1)$; gray marks the other branches. The level $t=0.05$ has two positive preimages. Only $v\approx3.1183$ belongs to the selected branch; $v\approx6.3817$ must be rejected. Open circles mark the tangencies.}

\clearpage
\subsection{The general-degree construction}\label{sec:generalcircle}
For each integer $p\ge3$, prepare $O_\Gamma=(h,j)$, $Z=(z_x,z_y)$, $k>0$ and $\rho>0$, using the rational formulas in Appendix~\ref{app:prep}. The circle passes through $B=(2,0)$ and is independent of $X$ and $s$. Put $\sigma=kX2^{p-3}/\rho$. The rail sequence uses
\[
 f_1=\rho,\quad f_i=\tfrac12\ (2\le i\le p-2),\quad f_{p-1}=\sigma.
\]
It saves $L_1=L(\rho s)$ and ends at $T=R(kt)$, because $\prod_i f_i=kX$. The same three final joins then give $R(sv)$.

\begin{theorem}[Counted finite-circle protocol]\label{thm:construction}
For the preparation in Appendix~\ref{app:prep}, assume $X,s>0$, $\tau_-<Xs^p<\tau_+$, and
\[
 s\ne1,\quad \sigma\ne1,\quad v\ne\rho,\quad
 b(v)=(2-h)2v-4j(\rho-v)\ne0.
\]
The selected branch is realized by $2p+2$ mobile joins using the prepared circle.
\end{theorem}
\begin{proof}
The power stage costs $1+2(p-1)=2p-1$ joins. To certify the remaining incidence, write
\[
 d=(2v,4(\rho-v)),\quad n=d\cdot d,\quad b=(B-O_\Gamma)\cdot d,
 \quad G=B-\frac{2b}{n}d.
\]
The pencil line through $B$ has direction $d$. Substituting $B+\xi d$ into the circle equation gives $\xi(\xi n+2b)=0$: besides $B$, the other intersection is exactly the displayed $G$. Thus $b\ne0$ keeps the two points distinct. The same line meets the top at $\hub{\rho/v}$ when $v\ne\rho$.

It remains to show that the projector $Z$ and the residual point $T$ select this $G$. Three points are collinear exactly when the determinant of their displacement vectors vanishes. Multiplying that determinant by $n$ gives
\begin{align*}
 I={}&(2-z_x)(-2bd_y-z_yn)\\
 &-[4(1-kt)-z_y][(2-z_x)n-2bd_x].
\end{align*}
The prepared parameters satisfy $I=\lambda_p[N(v)-tD(v)]$, with $\lambda_p\ne0$. This identity converts the geometric incidence into exactly the inverse model, rather than an approximation to it. Since the selected root has $N(v)-tD(v)=0$, it places $G$ on $ZT$. Projecting $G$ through $B$, then $L_1$ through the resulting center, gives $R(sv)$ by~\eqref{eq:railops}. These are three joins. The branch arc is parameterized by $v_-<v<v_+$ and can be marked during preparation using its endpoints and $G(1)$. The listed hypotheses exclude an infinite center, a coincident $B,G$, and tangency; the other rail intersections are transverse.
\end{proof}
The exclusions concern this finite chart. In particular, $s=1$ sends the initial center $V$ to infinity even though the scalar correction can remain regular. A positive rescaling $\widetilde X=c^pX$, $\widetilde s=s/c$ preserves $t$ and conjugates every correction. A dyadic choice can keep the limiting state away from one, controlling the distant initial center near convergence. The circle-specific exclusions still require their own branch and chart checks.

The result is a concrete upper bound on mobile work, not merely an existence theorem for straightedge construction. Its advantage over the rational fifth-order baseline will be quantified after the other devices have been introduced.

\clearpage
\subsection{Strict error contraction on the whole real branch}\label{sec:branch}
The real band is a convergence region, resolving the V20 conjecture. Here ``contraction'' means strict error reduction relative to the root, not a pairwise Lipschitz bound.

\begin{theorem}[Whole-branch contraction]\label{thm:branch}
Let $p\ge3$ be an integer and $\tau_-<t<\tau_+$, $t\ne1$. For the descending inverse branch $v=v(t)$,
\[
 \left|\frac{\log t}{p}+\log v\right|<\frac{|\log t|}{p}.
\]
Consequently the scalar circle iteration stays inside this residual band and converges to $X^{-1/p}$ from every start in it. Its local order is five.
\end{theorem}
\begin{proof}
\textbf{Reduce the error comparison to a sign.} Parameterize the branch by $v$, so $t=N(v)/D(v)$. Put
\[
 h(v)=p\log v+2\log\frac{N(v)}{D(v)},\qquad
 \mathcal H_p(v)=2\alpha_p\gamma_p(v-1)^4-N(v)D(v).
\]
The difference of squared errors is $\log v\,h(v)/p$. It suffices to prove $h(v)<0$ for $1<v<v_+$; reciprocity gives $h(1/v)=-h(v)$ and handles the other half. Equivalently, for integer $p$ the desired sign is $(v-1)[D(v)^2-v^pN(v)^2]>0$.

\textbf{Control the shape of $h$.} The certificate~\eqref{eq:cert} gives
\[
 h'(v)=\frac{p\mathcal H_p(v)}{vN(v)D(v)}.
\]
Write $v=1+u$. The useful simplification is
\[
 \frac{\mathcal H_p(1+u)}{u^2}=\alpha_p\gamma_pu^2+12(p^2-4)u+12(p^2-16)-\frac{288}{u}-\frac{144}{u^2}.
\]
For $u>0$, its derivative is strictly positive. Since $\mathcal H_p(1)=-144$, $h$ first decreases and can then increase only once. Thus $h(1)=0$ and the single endpoint inequality $h(v_+)<0$ imply $h(v)<0$ throughout $(1,v_+)$. Section~\ref{sec:endpoint} proves the endpoint inequality for every integer degree.

\textbf{Pass from one step to convergence.} Let $z=\log t$ and $F(z)=\log[tv(t)^p]$. The inverse branch is continuous in the open band, $F(0)=0$, and $|F(z)|<|z|$ for $z\ne0$. Since $\tau_-\tau_+=1$, an orbit starting at $z_0$ remains in the compact interval $[-|z_0|,|z_0|]$ strictly inside the band. Its absolute errors decrease to a limit $\ell$. If $\ell>0$, a convergent subsequence has limit $z_*$ with $|z_*|=\ell$; continuity would give both $|F(z_*)|=\ell$ and $|F(z_*)|<\ell$. Hence $\ell=0$. The local order follows from~\eqref{eq:circleorder}.
\end{proof}
The first, strict-error assertion is \texttt{whole\_branch\_contraction} in \texttt{RectangleFixedCircleBranch}; its 33-declaration audit checks the sign and endpoint argument. The invariant-band convergence and inverse-branch assembly above are written corollaries. This distinction avoids calling the entire analytic/geometric program a single formal theorem.

\clearpage
\subsection{An endpoint certificate valid at every degree}\label{sec:endpoint}
Let $r=\sqrt{12(p^2-1)}$, $u_+=(6+r)/(p^2-4)$ and $v_+=1+u_+$. At this branch endpoint,
\[
 (p^2-4)u_+^2=12(u_++1),\qquad \frac{N(v_+)}{D(v_+)}=\tau_-.
\]
Therefore $h(v_+)<0$ is exactly $v_+^p<\tau_+^2$. The quadratic relation reduces the polynomials to linear expressions:
\begin{align*}
 N(1+u_+)&=\frac{6[2(2p+1)-(p^2-2p-2)u_+]}{p+2},\\
 D(1+u_+)&=\frac{6[2(2p-1)+(p^2+2p-2)u_+]}{p-2}.
\end{align*}
For $3\le p\le9$, rational brackets for $r$ give exact rational upper bounds on $v_+^pN(v_+)^2/D(v_+)^2$ below one:
\begin{center}\small
\begin{tabular}{@{}crrrrrrr@{}}\toprule
$p$ & 3 & 4 & 5 & 6 & 7 & 8 & 9\\\midrule
Lower $r$ & 9.79 & 13.41 & 16.97 & 20.49 & 24 & 27.49 & 30.98\\
Upper $r$ & 9.80 & 13.42 & 16.98 & 20.50 & 24.01 & 27.50 & 30.99\\
Ratio $<$ & .130 & .147 & .154 & .157 & .159 & .161 & .162\\\bottomrule
\end{tabular}
\end{center}
The displayed decimal endpoints are exact rationals; the ratio bounds are rounded upward. The checked substitution uses the upper bound for $v_+$, and the lower bound for $u_+$ in both reduced expressions, bounding $N$ upward and $D$ downward.

For $p\ge10$, $3.46p-1\le r\le3.47p$ implies $pu_+\le5$ and $D(v_+)\ge12.3N(v_+)$. For the latter inequality, clearing positive denominators leaves the certificate
\[
 409q^4+11314q^3+140038q^2+1027656q+3519360>0,\qquad q=p-10\ge0.
\]
Hence $v_+^p\le e^{pu_+}\le e^5<148.42<12.3^2\le[D(v_+)/N(v_+)]^2$. The finitely many rational certificates and this exponential estimate cover all integer $p\ge3$.

\fig{branch_contraction}{1}{Sampled contraction ratios and endpoint margins illustrate the proved inequalities. The left curves use $p=3,5,17$; the horizontal line is the strict-error threshold one. The right curve samples degrees 3--400. Its limiting margin is $2\log(7+4\sqrt3)-2\sqrt3\approx1.80373$. No uniform numerical bound such as $0.36$ is asserted as a theorem.}
The whole-band result still requires real branch membership. The next constructions provide globally defined alternatives; Section~\ref{sec:safeguard} uses AD to enter the band and to handle excluded finite charts.

\clearpage
\section{Four accelerated Pandrosion constructions}\label{sec:supports}
The fixed circle solves for a correction externally. A second architecture keeps the power endpoint $E=R(s^p)$ and changes the interception support inside the rectangle. Prepare $M=(0,H(1-X^{-1}))$. The slope of $ME$ is $m=H(1-t)/(WX)$. A line through $A$ of slope $k_s=H\kappa/(WX)$ meets the parallel to $ME$ through $P=R(s)$ at $T_A$. The horizontal through $T_A$ gives $P^+$, with
\begin{equation}\label{eq:common}
 s^+=s\frac{k_s}{k_s-m}=s\frac{\kappa}{\kappa-1+t}.
\end{equation}
Each support supplies a different $\kappa$. The common triangle stage uses horizontals through $B_j=(Ws^j,H(1-s^j))$ on $OB$ and $L_j=L(s^j)$. Parallels to $L_1B$ carry $B_j$ to $L_{j+1}$; the $p$th horizontal gives $E$. This is the triangle protocol of Section~\ref{sec:rails}, not the scaled rail telescope. MA and its figure remain in Appendix~\ref{app:ma}.

\subsection{AK and centered AD: orders two and three}

\textbf{AK: a fixed line.} Prepare $K=(W(1-X/p),0)$. The support $AK$ has $\kappa=p$, so $\map{AK}(s)=ps/(p-1+t)$, Newton under inversion. It has order two and converges from every positive start.

\textbf{AD: a centered length transfer.} Prepare
\[
 O_{\mathrm{AD}}=\left(W-\frac{2W}{p-1},\ H-\frac{H(p+1)}{X(p-1)}\right).
\]
Transfer $AE=Hs^p$ downward from this center to $D_{\mathrm{AD}}$ and join $AD_{\mathrm{AD}}$. Then $\kappa=[p+1+(p-1)t]/2$, giving
\begin{equation}\label{eq:halley}
 \map{AD}(s)=s\frac{p+1+(p-1)t}{p-1+(p+1)t}.
\end{equation}
This Halley map has cubic order and global monotone convergence. Earlier geometric interpretations appear in \cite{gander1985,scavo1995}; here the rail encoding and readout are shared with AK.

\begin{figure}[H]\centering
\begin{minipage}[t]{.48\linewidth}\centering
\includegraphics[width=.88\linewidth]{figures/pandrosion_ak.pdf}
\caption{Pandrosion AK, order two. The fixed support $AK$ replaces MA while the triangle chain and green correction direction are retained. Here $p=3,X=2,s=3/4$.}
\end{minipage}\hfill
\begin{minipage}[t]{.48\linewidth}\centering
\includegraphics[width=.88\linewidth]{figures/pandrosion_ad.pdf}
\caption{Pandrosion AD, order three. Transfer the orange length $AE$ onto the prepared vertical through $O_{\rm AD}$; the lower arc point determines $AD_{\rm AD}$. Here $O_{\rm AD}=C$.}
\end{minipage}
\end{figure}
Both drawings retain the triangle stage, green transport and horizontal readout; only the support changes.
\clearpage
\subsection{Projective AD: order four}
\textbf{Two joins supply the support.} Prepare
\[
 \ell:\ y=H\left(1+\frac{5p-1}{(p+1)X}\right),\quad
 Z_{\mathrm{proj}}=\left(W\left(1+\frac{5p-1}{2p(2p-1)}\right),H\left(1+\frac{p+1}{(2p-1)X}\right)\right).
\]
Set $D_{\mathrm{proj}}=Z_{\mathrm{proj}}E\cap\ell$ and join $AD_{\mathrm{proj}}$. Substitution in~\eqref{eq:common} gives
\begin{equation}\label{eq:proj}
 \map{proj}(s)=s\frac{2p((2p-1)t+p+1)}{(p+1)t^2+2(2p-1)(p+1)t+(2p-1)(p-1)}.
\end{equation}
The direct-root correction is Pad\'e $[2/1]$; the reciprocal-state factor displayed here is $[1/2]$. Denote its positive denominator by $Q_{\mathrm{proj}}(t)$. Direct differentiation gives
\[
 \map{proj}'(s)=-\frac{2p(p^2-1)(2p-1)(t-1)^3}{Q_{\mathrm{proj}}(t)^2}.
\]
The map has its global maximum $a$ at $s=a$. Also $\map{proj}(s)/s-1=(1-t)((p+1)t+5p-1)/Q_{\mathrm{proj}}(t)$. Thus a first step lands at or below $a$, and subsequent steps increase to $a$. Integrating the derivative at $a$ gives exact order four. The archival companion gives further dynamics and the original triangle-chain realization \cite{geometryv20}.

\fig{pandrosion_projective}{.49}{Projective AD [2/1], order four, at the same cubic state. The orange line $Z_{\rm proj}E$ meets the prepared horizontal $\ell$ at $D_{\rm proj}$. The blue support $AD_{\rm proj}$ receives the green correction. The inset separates $P$, $T$ and $P^+$.}
\textbf{The fifth construction.} Displacing the arc center from the prepared vertical makes the height a genuine square root. The decentered AD support is developed next, because its global fifth-order behavior changes the choice of device rather than only its coefficients.

\clearpage
\subsection{A decentered arc with a global fifth-order map}
For $p>2$, define
\begin{equation}\label{eq:arc}
 q(t)=\sqrt{\frac{(p-2)(t^2+1)+(10p+4)t}{12p}},\quad g(t)=(p-1)q(t),\quad
 \map{arc}(s)=s\frac{1+g(t)}{t+g(t)}.
\end{equation}
For integer $p\ge3$, put $\eta=(p-1)\sqrt{(p-2)/(12p)}$, $\zeta=(5p+2)/(p-2)$, $d_0=W/\eta$ and $\delta=(H/X)\sqrt{\zeta^2-1}$. Prepare
\[
 G_0=(W,H+H\zeta/X),\quad
 O_{\mathrm{arc}}=(W-d_0+\delta,H-H/(X\eta)),\quad x=W-d_0.
\]
Draw the arc centered at $O_{\mathrm{arc}}$ with radius $EG_0=(H/X)(t+\zeta)$ and take the lower intersection $D_{\mathrm{arc}}$ with the prepared vertical. Its vertical distance below the center is
\[
 \sqrt{EG_0^2-\delta^2}=\frac HX\sqrt{t^2+2\zeta t+1}.
\]
The support $AD_{\mathrm{arc}}$ consequently gives $\kappa=1+g(t)$ in~\eqref{eq:common}. All preparation parameters use rational operations and square roots of positive quantities. The radicand stays positive for $t>0$.
\fig{decentered_arc}{1}{The decentered-arc device for $p=7,X=W=2,H=4,s=3/4$. The full circle geometry is shown separately from the local report, at independently chosen equal-axis scales. The displacement $\delta$ makes the vertical circle intersection perform a square root. The local panel retains the complete support and distinguishes the reported state.}
Unlike the inverse-circle correction, this map is defined for every positive residual and requires no choice between two inverse-model branches. It uses a mobile radius at every step. Unlike rational diagonal Pad\'e, it realizes the correction with a single radical rather than two rational factors. These are practical differences in geometric realization, even though all three maps have fifth local order.

Centered AD and decentered AD use the same mobile support operations: one arc and one join. Their fixed preparations and coordinate scales differ. This is the precise sense in which order five is obtained at the mobile support cost of order three.

\clearpage
\subsection{Global monotone convergence of the decentered map}
\begin{theorem}\label{thm:arc}
For real $p>2$, $X>0$ and every $s_0>0$, \eqref{eq:arc} converges to $a$ monotonically without crossing it. Its logarithmic error satisfies
\begin{equation}\label{eq:arcorder}
 e^+=\frac{(p-2)(p^2-1)(3p+1)}{1440}e^5+O(e^7).
\end{equation}
\end{theorem}
\begin{proof}
The radicand is positive and $q(1/t)=q(t)/t$, so the correction is reciprocal. Write
\[
 \mathcal E(t)=\frac{\log t}{p}+\log(1+g(t))-\log(t+g(t)).
\]
Define the auxiliary polynomials
\begin{align*}
 A_*(t)&=(p-1)(t^2+1)+(10p+2)t,\\
 B_*(t)&=(t+1)(6pt-(t-1)^2),\\
 K_*(t)&=(p^2-5p-2)(t-1)^2+12p(3p+1)t.
\end{align*}
Differentiation and factorization give
\begin{align}
 \frac{d\mathcal E}{d\log t}
 &=\frac{(p-2)(p-1)^2(qA_*-B_*)}{12p^2g(1+g)(t+g)},\label{eq:arcder}\\
 q^2A_*^2-B_*^2&=\frac{p+1}{12p}(t-1)^4K_* .\label{eq:arcfactor}
\end{align}
Here $A_*>0$. If $B_*\le0$, then $qA_*-B_*>0$. Otherwise $(t-1)^2<6pt$. When $p^2-5p-2\ge0$, $K_*>0$ directly; otherwise
\[
 K_*=6p^2(p+1)t+(p^2-5p-2)((t-1)^2-6pt)>0.
\]
Equation~\eqref{eq:arcfactor} now gives $qA_*>B_*$ for $t\ne1$. Therefore $\mathcal E$ is strictly increasing with $\mathcal E(1)=0$: the new logarithmic error has the old sign. Meanwhile $\Phi(t)-1=(1-t)/(t+g)$ moves the state toward the root. The orbit is monotone and bounded between $s_0$ and $a$. Its positive limit is a fixed point by continuity, and $\Phi(t)=1$ implies $t=1$.

Near one, rationalize $qA_*-B_*$ using~\eqref{eq:arcfactor}. Since $q=1$, $g=p-1$, $A_*=B_*=12p$ and $K_*=12p(3p+1)$ at $t=1$, this gives
\[
 \lim_{t\to1}\frac{d\mathcal E/d\log t}{(t-1)^4}
 =\frac{(p-2)(p^2-1)(3p+1)}{288p^5}.
\]
Substituting $t=e^{pe}$ and integrating yields~\eqref{eq:arcorder}; reciprocity makes the analytic logarithmic map odd.
\end{proof}
At $p=3$, the leading coefficient is $+1/18$, whereas the fixed-circle inverse has $-1/18$. The former stays on the starting side; the latter alternates near the root. At $p=2$, both scalar radical formulas reduce to the exact square-root correction, but the displayed decentered preparation must be replaced by a mean-proportional construction.

\clearpage
\section{What the shorter construction saves}\label{sec:cost}
A useful comparison holds the state, requested accuracy and counting convention fixed. The rational competitors can be built on the same rail telescope. A prepared projector
\[
 Z(\mu,\nu)=\left(-\frac{W\mu}{1-\mu},H\left(1-\frac{\nu}{1-\mu}\right)\right)
\]
sends $Q=R(q_p)$ to $\hub{f(q_p)}$, with $f(q)=\mu+\nu/q$. Two such projections and two rail joins multiply a stored value by $f/g$: a four-join module.

For Halley choose $f=(p-1)/(4p)+(p+1)/(4pt)$ and $g=(p+1)/(4p)+(p-1)/(4pt)$. Since $t=X2^{p-1}q_p$, each is of the required projector form. The cost is $2p-1+4=2p+3$ joins.

For rational diagonal Pad\'e set $a_5=(p-1)(2p-1)$, $b_5=8p^2-2$, $c_5=(p+1)(2p+1)$. The correction is
\[
 \Phi_{\mathrm{rat}}(t)=\frac{a_5t^2+b_5t+c_5}{c_5t^2+b_5t+a_5}
 =\prod_{r\in\{r_+,r_-\}}\frac{t+r}{rt+1},\quad
 r_\pm=\frac{b_5\pm\sqrt{12p^2(4p^2-1)}}{2a_5}.
\]
Each factor uses the four-join module with $f_r=[1+r/t]/[2(1+r)]$ and $g_r=[r+1/t]/[2(1+r)]$. Both reuse $Q$, giving $2p+7$ joins. Its scalar map is globally monotone: its derivative is $a_5c_5(t-1)^4/(c_5t^2+b_5t+a_5)^2$, and numerator minus denominator is $6p(1-t)(1+t)$.

\begin{table}[H]\centering\small
\caption{Generic mobile protocol costs; fixed preparation and final inversion are excluded. A parallel macro is $2J+3C_{\mathrm{arc}}$. The scope column describes the exact scalar map.}
\begin{tabularx}{\linewidth}{@{}Xccccc@{}}\toprule
Construction & Order & $J$ & $P_{\parallel}$ & $C_{\mathrm{arc}}$ & Scope\\\midrule
MA rectangle & 1 & 2 & $2p+1$ & 0 & $X>1$\\
AK rectangle & 2 & 2 & $2p+1$ & 0 & Global\\
AD rectangle & 3 & 3 & $2p+1$ & 1 & Global\\
Projective AD & 4 & 4 & $2p+1$ & 0 & Global\\
Decentered AD & 5 & 3 & $2p+1$ & 1 & Global\\
Pencil Halley & 3 & $2p+3$ & 0 & 0 & Global\\
Pencil rational Pad\'e & 5 & $2p+7$ & 0 & 0 & Global\\
Fixed-circle inverse & 5 & $2p+2$ & 0 & 0 & Real band\\
Circle--AD protocols & 5 & \multicolumn{3}{c}{Section~\ref{sec:safeguard}} & Global\\\bottomrule
\end{tabularx}
\end{table}
The fixed circle saves five joins per regular update relative to this rational fifth-order realization. A reproducible cubic example makes both the gain and the price of control visible.
\begin{table}[H]\centering\small
\caption{Steps / mobile joins to obtain $|s/a-1|\le10^{-d}$ from $p=3,X=2,s_0=3/4$, computed at 250 digits. The last two rows use binary residual caching and add 12 initial joins; every circle trial in this example is regular and accepted.}
\begin{tabularx}{\linewidth}{@{}Xcccc@{}}\toprule
Protocol & $d=6$ & $d=12$ & $d=30$ & $d=60$\\\midrule
Pencil Halley & 2 / 18 & 3 / 27 & 3 / 27 & 4 / 36\\
Pencil rational Pad\'e & 1 / 13 & 2 / 26 & 2 / 26 & 3 / 39\\
Circle, linear power & 1 / 8 & 2 / 16 & 2 / 16 & 3 / 24\\
Circle, branch entry & 1 / 18 & 2 / 36 & 2 / 36 & 3 / 54\\
Circle, residual gate & 1 / 28 & 2 / 56 & 2 / 56 & 3 / 84\\
\bottomrule\end{tabularx}
\end{table}
These are geometric work counts under the stated convention, not circuit time or energy measurements. The next section develops branch entry; the optional residual gate is documented in Appendix~\ref{app:gate}.


\clearpage
\section{From a convergence band to all positive starts}\label{sec:safeguard}
Use a regular descending circle readout whenever it exists; otherwise use centered AD. The branch theorem removes any need to test its improvement. The older, independently formalized residual gate is retained in Appendix~\ref{app:gate}.

\subsection{AD entry and geometric exceptions}
Cache $E=R(s^p)$ and $R(t)$. Stop at $t=1$; otherwise apply the chosen step and refresh its power and residual.

\begin{proposition}[Global branch-driven protocol]\label{thm:entry}
For every integer $p\ge3$, $X>0$ and $s_0>0$, this exact geometric protocol converges to $a=X^{-1/p}$. With the preparation specified below, it eventually uses the circle step unless already at the root, and has local order five.
\end{proposition}
\begin{proof}
Outside the band, AD converges globally and therefore reaches the open band in finitely many steps. Inside, both maps strictly reduce $|\log t|$ and preserve the band. On the compact log-residual interval determined by entry, both scalar post-error functions are continuous. If the errors tended to $\ell>0$, the larger post-error would be uniformly below the current error on $\ell\le|\log t|\le|\log t_{\rm entry}|$. Every late step would decrease it by a fixed positive amount, a contradiction. Thus chart-driven switching also converges.

The finite chart is regular near the root. For the compact preparation ($p\ne4$), $\rho=1/2$ and
\[
 b(1)=\frac{48p(p+4)(p+1)}{(p-1)\Lambda_p}>0;
\]
for $p=4$, use the separate chart in Appendix~\ref{app:prep}, where direct substitution also gives $\rho\ne1$ and $b(1)\ne0$. At the root $\Delta_p(1)=36p^2>0$. Continuity therefore preserves the transverse crossing, distinct $B,G$ and finite top intersection nearby; the eventual update is the order-five circle map.
\end{proof}
Lean checks the branch inequality and chart certificates separately; the switching assembly above is a written proof.

\subsection{Geometric arithmetic and complete trace counts}
For positive $a,b$, reuse $R(a)$ when $b=1$. Otherwise one join gives $V=CR(b)\cap OA=\hub{1/b}$, then four joins through $\hub{1/2},V,\hub{1/2},\hub{1/4}$ successively give $L(a/2),R(ab/2),L(ab/4),R(ab)$. This five-join product is finite, including its exact unit shortcut. The two centers $\hub{d/[2(1+d)]}$ and $\hub{1/[2(1+d)]}$ multiply a rail value by a fixed $d>0$ in two joins.

With $m=\lfloor\log_2p\rfloor+\operatorname{popcount}(p)-1$, binary powering and the fixed scale $X$ cost $5m+2$ joins. Each power makes $m$ unit tests; branch and root-stop decisions are additional, uncharged decisions.

\begin{table}[H]\centering\small
\caption{Mobile upper bounds, including the next cache. Fixed preparation, existing intersections and decisions are excluded.}
\begin{tabularx}{\linewidth}{@{}Xccc@{}}\toprule
Operation & $J$ & $P_{\parallel}$ & $C_{\mathrm{arc}}$\\\midrule
Initial power and residual cache & $5m+2$ & 0 & 0\\
Branch-driven regular circle step & $5m+8$ & 0 & 0\\
Branch-driven fallback after a full attempted readout & $5m+10$ & 2 & 1\\
\bottomrule
\end{tabularx}
\end{table}
A trial costs six joins: two for $R(kt)$, one for $L(\rho s)$ and three for the circle readout. AD adds $2J+2P_{\parallel}+C_{\rm arc}$. The cubic costs $12J$ initially, $18J$ per regular step and $31$ expanded traces for full fallback. At $p=10^6$, $m=25$ gives $133J$ per regular step. These are instrumented protocol bounds, not a global eight-line claim.

\clearpage
\section{Binary geometry and the path to an analog calculator}\label{sec:binary}
The linear telescope is attractive for a few degrees. At large $p$, binary powering gives a smaller arithmetic architecture. Define a top encoding $Z_\times(q)=(W+Hq,H)$ and the fixed unit $U_\times=(W+H,H)$. A parallel to $U_\times B$ copies $R(a)$ to $Z_\times(a)$. The parallel through $Z_\times(a)$ to $U_\times R(b)$ then meets the right rail at $R(ab)$: its slope is $b$ and its horizontal displacement is $-Ha$.

Squaring the accumulator and multiplying by $s$ at each binary one produces $E=R(s^p)$. With $L=\lfloor\log_2p\rfloor$ and $h=\operatorname{popcount}(p)$, there are $m=L+h-1$ multiplications. Reusing the top copy of $s$, the power stage costs $mJ+(m+L)P_{\parallel}$. It feeds any of the four rectangle supports unchanged.
\fig{fast_geometry}{1}{The fast-power construction and its AD readout, adapted from the V30 companion. The case $p=13=(1101)_2$ uses five products along exponents $1,2,3,6,12,13$. A product is an actual pair of line incidences, not a numerically inserted power. The output $E$ enters the same moving support and horizontal readout.}
The construction is an operation graph for a circuit: it identifies which signals must exist before a cell can act, and which intermediate values must survive a later correction. The implementation obligations are concrete:
\begin{center}\small
\begin{tabularx}{\linewidth}{@{}lXX@{}}\toprule
Geometric object & Electrical role & Property to measure\\\midrule
Rail coordinate & Scaled voltage or current & Dynamic range and offset\\
Product incidence & Multiplier and weighted sums & Gain error and settling\\
Circle crossing & Quadratic solve plus branch choice & Stability and branch margin\\
Retained point & Sample-and-hold memory & Leakage and sampling error\\\bottomrule
\end{tabularx}\end{center}
We now separate the existing AD circuit from the quadratic cell suggested by the fixed-circle construction.

\clearpage
\subsection{From an incidence equation to an analog cell}
\textbf{What has been implemented.} The V30 prototype realizes centered AD, of order three. A digitally prepared scale $c$ makes $Y\approx Xc^p\in[1,2]$, with $s=cu$. Storing $u\approx1$ directly would give sensitivity $\partial\log(Yu^p)/\partial u=p/u$. Instead store $q=p(u-1)$: the sensitivity becomes $1/u$. Centered binary products evaluate the residual, followed by
\[
 q^+=q+\frac{2(1+q/p)(1-t)}{1+t+(t-1)/p},\qquad
 \widehat r=\frac1{c(1+q/p)}.
\]
The V30 companion gives the exact centered product recurrences, netlists and calibration protocol \cite{analogv30}. There are 25 power products at $p=10^6$, and at most 37 over the programmed range $3\le p\le10^6$.
\fig{analog_circuit}{.73}{Existing AD behavioral topology from V30. State memory, loading and switches are explicit; diamonds are behavioral arithmetic sources. The drawing identifies where state must settle and be retained. It is the implemented order-three baseline, not an electrical implementation of the circle solver.}

\textbf{What the fixed circle specifies for a new cell.} Its geometric incidence is an exact circuit specification: solve $N(v)-tD(v)=0$ on the branch through $v=1$, then form $s^+=sv$. A circle on paper does not by itself supply an electrical component. A candidate realization would use multipliers, weighted sums and feedback to enforce this quadratic equation, with branch selection and AD fallback.

As with AD, the useful electrical variable at large $p$ is centered: put $z=p(v-1)$. Dividing the incidence equation by 12 gives the exact relation
\begin{align*}
 0={}&(1-t)-\frac{(1-2/p)+(1+2/p)t}{2}\,z\\
 &+\frac{(1-3/p+2/p^2)-t(1+3/p+2/p^2)}{12}\,z^2 .
\end{align*}
At $(t,z)=(1,0)$ the derivative of the left side with respect to $z$ is $-1$, independently of $p$. This removes the trivial large-degree scaling from the local quadratic cell. At a branch endpoint that derivative vanishes: an electrical design must restrict its operating margin or fall back before the tangency. Algebraic centering alone does not control noise or feedback settling.

\clearpage
\subsection{What the prototype teaches about the construction}
The existing campaigns distinguish programmability from electrical accuracy. Arithmetic checks cover 8,243 $(p,X)$ pairs. The calibrated behavioral SPICE study meets one ppm decoded-root error in 64 cases, with worst error 0.279 ppm at 23.985 ms. A 4.785 ms policy meets one ppm on 18 new cases. These results expose a precision--time tradeoff in a behavioral model; they are not silicon measurements.

Preparation matters: at $p=10^6,X=500000$, the digitally prepared $1/c$ is already within about $2.477\times10^{-7}$ relative error of the root. The analog loop refines it, and its recorded readout time excludes preparation, programming and conversion.

The next experiment should retain the same power module and memories, replace AD by the centered quadratic cell, and measure wrong-branch events, settling and root error under identical loading and noise. Order five may save updates, but each nonlinear solve has a cost. No electrical advantage has yet been established.

\section{Evidence and reproducibility}\label{sec:scope}
The main proof boundary is explicit: Lean establishes the whole-branch strict-error inequality; the invariant-band, inverse-branch/local-order and branch-driven convergence assemblies are written arguments. Appendix~\ref{app:gate} has a separate formal global switching theorem. The decentered map's global radical proof is written. The source map below identifies the checked components \cite{lean4,mathlib}.
\begin{center}\small
\begin{tabularx}{\linewidth}{@{}p{.34\linewidth}X@{}}\toprule
Lean module & Checked component\\\midrule
\texttt{RectangleProjective} & Incidences and order-four dynamics.\\
\texttt{RectanglePencils*} & Rail/projector identities, dynamics and orders.\\
\texttt{RectangleFixedCircleBranch} & Strict branch contraction and endpoint bounds; 33 declarations audited.\\
\texttt{RectangleFixedCircle} & Quadratic identities, charts, preparation, readout and a parameterized local limit.\\
\texttt{RectangleDecenteredArc} & Certificates for the written radical proof.\\
\texttt{RectangleCircleSafeguard} & Products, powers, residual gate, global switching and local transfer; 41 declarations audited.\\
\texttt{RectangleFastPower} & Product incidence and binary recursion.\\
\texttt{AnalogFastAD} & Six centered-arithmetic identities.\\\bottomrule
\end{tabularx}\end{center}
Both audits allow only \texttt{propext}, \texttt{Classical.choice} and \texttt{Quot.sound}. Physical trace counts are instrumented inventories, not a Lean cost semantics. Neither the proofs nor sampled contraction ratios certify uniform drawing size, browser accuracy or hardware performance.

Independent intersections check 162 circle configurations at 250 digits, 108 arc configurations at 160 digits, and two 144-trajectory controller campaigns at 180 digits. These finite tests supplement the proofs. The \href{https://github.com/ivan-fr/pandrosion/blob/main/paper/reciprocal_geometry_v22/README.md}{repository guide} links all sources, audits, data and reproduction commands. CI rebuilds the PDF; the \href{https://ivan-fr.github.io/pandrosion/}{gallery} provides interactive views. The repository includes the MIT license and citation metadata; V20, V21 and analog V30 remain archived.

The numerical families are classical \cite{dlmf38,laszkiewicz2009,baker1996}; priority and trace optimality of these realizations remain open. Iterative approximation is distinct from exact finite construction.

\section{Conclusion}
The prepared circle gives fifth-order convergence on its real band in $2p+2$ joins; AD entry and chart fallback extend it globally at an explicit extra cost. The centered quadratic relation now supplies a testable analog-cell design. Its error and settling budget must be measured against the existing AD baseline.

\clearpage
\appendix
\section{Rational preparation for every integer degree}\label{app:prep}
This appendix makes the general-circle protocol reproducible without searching through an earlier manuscript. Write $\nu=p-3$ and use $\alpha_p,\gamma_p$ from~\eqref{eq:model}. Define
\begin{align*}
 \mathcal D_p&=625\nu^5+10525\nu^4+56797\nu^3+88639\nu^2+14758\nu+664,\\
 \mathcal K_p&=625\nu^5+1975\nu^4+4147\nu^3+2905\nu^2+448\nu+340,\\
 \chi_p&=288p(p-2)(25p^2-244)/\mathcal D_p,\\
 u_p&=5-4\alpha_p/(9\gamma_p),\qquad w_p=8(5-2p)/(9(p-1)).
\end{align*}
Take
\begin{align*}
 \rho&=\tfrac13,& k&=\mathcal K_p/\mathcal D_p,& j&=-3\chi_pw_p/4,\\
 h&=2+2j-\chi_pu_p/2,&z_x&=2-\chi_p,&
 z_y&=\frac{384p(25p^3-50p^2-415p+974)}{\mathcal D_p}.
\end{align*}
Then the incidence certificate in Theorem~\ref{thm:construction} has multiplier $\lambda_p=-64k\chi_p/(9\gamma_p)$. The positive coefficients imply $\mathcal D_p,\mathcal K_p>0$ for $\nu\ge0$; $25p^2\ne244$ for integer $p\ge3$ makes $\chi_p,\lambda_p,j$ nonzero. All parameters are rational. The circle is drawn with center $(h,j)$ through $B$. This uniform chart is checked symbolically; the Lean preparation certificates cover the compact chart below and the separate degree-four specialization.
At $p=4$ the formulas give the finite values
\[
 k=145/2389,\quad O_\Gamma=(13502/7167,3328/2389),\quad
 Z=(-214/2389,2432/2389).
\]
The cubic worked example uses a more compact chart. For $p\ne4$, put $\Lambda_p=p^3+12p^2+39p-26$ and take
\begin{align*}
 \rho&=\tfrac12,& k&=\frac{(p-2)(p^2-4p+13)}{\Lambda_p},\\
 h&=2-\frac{72p(p+4)}{(p-1)\Lambda_p},&
 j&=\frac{24p(p-2)(p+4)}{(p-1)\Lambda_p},\\
 z_x&=2-\frac{24p(p+4)(p-2)}{(p-4)\Lambda_p},&
 z_y&=\frac{24p(p^2+2p-26)}{(p-4)\Lambda_p}.
\end{align*}
Here $\lambda_p=-384p(p-2)(p+4)(p^2-4p+13)/[(p-4)(p-1)\Lambda_p^2]$. Substitution checks both charts. For numerical evaluation the branch can be written without the removable cancellation at large $t$:
\[
 v(t)=\begin{cases}
 \dfrac{\alpha_p-t\gamma_p}{\beta_p(1-t)+\sqrt{\Delta_p(t)}},&t\le1,\\[5pt]
 \dfrac{\sqrt{\Delta_p(t)}-\beta_p(1-t)}{t\alpha_p-\gamma_p},&t>1.
 \end{cases}
\]

\section{The linear MA reference}\label{app:ma}
MA provides a baseline, not an additional accelerated method. Write its lower-right endpoint as $A_{\rm lin}=B$. The parallel to $ME$ through $P$ meets $MB$; a horizontal returns the next state. For $s\ne1$,
\[
 \map{MA}(s)=1-\frac{X-1}{X S_p(s)},\qquad S_p(s)=\sum_{j=0}^{p-1}s^j.
\]
For $X>1$ this scalar map is increasing and $\map{MA}(s)-s=(1-Xs^p)/(XS_p(s))$. It therefore converges monotonically to $a$ from every positive start, with the removable scalar extension at $s=1$. Its derivative at $a$ is $1-pa^{p-1}/S_p(a)\in(0,1)$: convergence is linear.
\fig{pandrosion_ma}{.34}{Pandrosion MA, the linear reference construction, at $p=3,X=2,s=3/4$. Thin lines construct the power by similar triangles. The green parallel meets the blue support $MA_{\rm lin}$; the horizontal gives $P^+$. The historical endpoint $A_{\rm lin}$ is the lower-right vertex $B$ in this paper. At $s=1$ the geometric intersection is coincident despite the removable scalar extension.}

The notation $S_p(s)$ is reserved for this geometric sum; the sign polynomial in the branch proof is $\mathcal H_p(v)$.

\section{The independent residual-gate controller}\label{app:gate}
This independently formalized V21 controller is retained for comparison. It is not required by the branch-driven protocol.

For a positive circle candidate $c$, construct $u=Xc^p$ and accept precisely when
\begin{equation}\label{eq:gate}
 (t\le1\ \wedge\ t\le u^2\ \wedge\ tu^2\le1)
 \quad\hbox{or}\quad
 (1\le t\ \wedge\ u^2\le t\ \wedge\ 1\le tu^2).
\end{equation}
Otherwise use AD. Three rail comparisons implement this condition, equivalent to $|\log u|\le|\log t|/2$. The errors decrease; if their limit were positive, eventually no halving step would be possible, leaving only AD and forcing the limit to zero. This is the V21 global switching theorem, including the incidence-based gate, formalized in \texttt{RectangleCircleSafeguard}. Strict branch contraction alone does not assert the stronger half-error inequality.

From the cached current state, a circle trial uses six joins and the candidate power/residual uses $5m+2$. The gate's two products add at most ten joins and three ordered comparisons. Acceptance thus costs $5m+18$ joins including the next cache. Full rejection adds AD ($2J+2P_{\parallel}+C_{\rm arc}$) and a fresh $5m+2$ cache, for a total of $(10m+22)J+2P_{\parallel}+C_{\rm arc}$. The common initialization costs $5m+2$ joins.

\begin{center}\small
\begin{tabular}{@{}lrr@{}}\toprule
Cubic protocol ($m=2$) & Regular / accepted & Full fallback\\\midrule
Branch-driven & $18J$ & $20J+2P_{\parallel}+C_{\rm arc}$\\
Residual gate & $28J$ & $42J+2P_{\parallel}+C_{\rm arc}$\\\bottomrule
\end{tabular}
\end{center}
The expanded fallback bounds are respectively 31 and 53 traces. At $p=10^6$, regular steps cost 133 versus 143 joins. A completed gate adds two product-identity tests to the $m$ tests in each power evaluation. All these are protocol upper bounds; fixed preparation and decisions are not charged as traces.

\fig{circle_safeguard}{.82}{The five-join product and the optional V21 residual gate. The branch-driven V22 protocol reuses the product for powers and scales, but omits the two gate products and its three ordered comparisons. The geometric points and labels are unchanged from the independently checked V21 figure.}
Its separate 144-trajectory campaign and 41-declaration axiom audit are reproduced in the repository.
\begingroup\footnotesize\setlength{\bibsep}{1pt}
\bibliographystyle{unsrtnat}\bibliography{references}\endgroup
\end{document}
